\textbf{Phần 1. Trắc nghiệm lựa chọn}
\setcounterpageref{ex}{0}

\begin{ex} %[4D10V53]
Phân số chỉ phần đã tô màu gồm 3 phần bằng nhau trong tổng số 4 phần bằng nhau của hình tròn là
\choice
{$\dfrac{4}{3}$}
{$\dfrac{1}{4}$}
{\True $\dfrac{3}{4}$}
{$\dfrac{3}{7}$}
\loigiai{
Phân số chỉ phần đã tô màu là $\dfrac{3}{4}$.\\
\textbf{Đáp án của câu hỏi là C}
}
\end{ex}

\begin{ex} %[4D10V53]
Trong phân số $\dfrac{5}{8}$, tử số là
\choice
{\True $5$}
{$8$}
{$13$}
{$3$}
\loigiai{
Tử số là số viết trên gạch ngang: 5.\\
\textbf{Đáp án của câu hỏi là A}
}
\end{ex}

\begin{ex} %[4D10V53]
Trong phân số $\dfrac{7}{12}$, mẫu số là
\choice
{$7$}
{$19$}
{$5$}
{\True $12$}
\loigiai{
Mẫu số là số viết dưới gạch ngang: 12.\\
\textbf{Đáp án của câu hỏi là D}
}
\end{ex}

\begin{ex} %[4D10V53]
Phân số \textit{Bảy phần mười một} được viết là
\choice
{$\dfrac{11}{7}$}
{\True $\dfrac{7}{11}$}
{$7{,}11$}
{$11{,}7$}
\loigiai{
Viết là $\dfrac{7}{11}$.\\
\textbf{Đáp án của câu hỏi là B}
}
\end{ex}

\begin{ex} %[4D10V53]
Mẫu số của một phân số luôn phải thỏa mãn điều kiện nào?
\choice
{Là số tự nhiên lớn hơn 10}
{Bằng tử số}
{\True Là số tự nhiên khác 0}
{Nhỏ hơn tử số}
\loigiai{
Mẫu số phải là số tự nhiên khác 0.\\
\textbf{Đáp án của câu hỏi là C}
}
\end{ex}

\begin{ex} %[4D10V53]
Một cái bánh pizza được chia đều thành 8 miếng bằng nhau. Bạn An đã ăn 3 miếng. Phân số chỉ phần bánh còn lại là
\choice
{\True $\dfrac{5}{8}$}
{$\dfrac{3}{8}$}
{$\dfrac{8}{3}$}
{$\dfrac{5}{3}$}
\loigiai{
Phần còn lại: $8 - 3 = 5$ miếng $\Rightarrow \dfrac{5}{8}$ cái bánh.\\
\textbf{Đáp án của câu hỏi là A}
}
\end{ex}

\begin{ex} %[4D10V53]
Một hộp có 12 quả bóng gồm 5 bóng đỏ và 7 bóng xanh. Phân số chỉ số bóng đỏ so với tổng số bóng trong hộp là
\choice
{$\dfrac{7}{12}$}
{$\dfrac{5}{7}$}
{$\dfrac{7}{5}$}
{\True $\dfrac{5}{12}$}
\loigiai{
Phân số là $\dfrac{5}{12}$.\\
\textbf{Đáp án của câu hỏi là D}
}
\end{ex}

\begin{ex} %[4D10V53]
Một lớp học có 35 học sinh, trong đó có 18 học sinh nữ. Phân số chỉ số học sinh nam so với cả lớp là
\choice
{$\dfrac{18}{35}$}
{\True $\dfrac{17}{35}$}
{$\dfrac{17}{18}$}
{$\dfrac{18}{17}$}
\loigiai{
Số nam: $35 - 18 = 17$ học sinh $\Rightarrow \dfrac{17}{35}$.\\
\textbf{Đáp án của câu hỏi là B}
}
\end{ex}

\begin{ex} %[4D10V53]
Trên một đoạn thẳng $AB$ dài $10\text{ cm}$, điểm $M$ chia đoạn thẳng $AB$ thành 5 phần bằng nhau. Điểm $C$ nằm ở vạch chia thứ 2 tính từ $A$. Phân số chỉ độ dài đoạn $AC$ so với $AB$ là
\choice
{$\dfrac{3}{5}$}
{$\dfrac{2}{10}$}
{\True $\dfrac{2}{5}$}
{$\dfrac{1}{5}$}
\loigiai{
$AC$ chiếm 2 phần trong tổng 5 phần $\Rightarrow \dfrac{2}{5}$.\\
\textbf{Đáp án của câu hỏi là C}
}
\end{ex}

\begin{ex} %[4D10V53]
Viết phân số có tử số là số chẵn bé nhất có 2 chữ số và mẫu số là số lớn nhất có 2 chữ số.
\choice
{\True $\dfrac{10}{99}$}
{$\dfrac{12}{99}$}
{$\dfrac{10}{98}$}
{$\dfrac{2}{99}$}
\loigiai{
Tử số là 10, mẫu số là 99 $\Rightarrow \dfrac{10}{99}$.\\
\textbf{Đáp án của câu hỏi là A}
}
\end{ex}

\begin{ex} %[4D10V53]
Một can đựng đầy $10\text{ l}$ nước. Người ta đã dùng $7\text{ l}$ nước. Phân số chỉ số lít nước đã dùng là
\choice
{$\dfrac{3}{10}$}
{$\dfrac{7}{3}$}
{$\dfrac{10}{7}$}
{\True $\dfrac{7}{10}$}
\loigiai{
Đã dùng $\dfrac{7}{10}$ can nước.\\
\textbf{Đáp án của câu hỏi là D}
}
\end{ex}

\begin{ex} %[4D10V53]
Có bao nhiêu phân số có mẫu số bằng 7 mà tử số là các số tự nhiên khác 0 và nhỏ hơn mẫu số?
\choice
{$5$ phân số}
{\True $6$ phân số}
{$7$ phân số}
{$8$ phân số}
\loigiai{
Các tử số là $1, 2, 3, 4, 5, 6$ (có 6 phân số: $\dfrac{1}{7}, \dots, \dfrac{6}{7}$).\\
\textbf{Đáp án của câu hỏi là B}
}
\end{ex}

\textbf{Phần 2. Trắc nghiệm đúng sai}
\setcounterpageref{ex}{0}

\begin{ex} %[4D10V53]
Xét tính đúng sai về cấu tạo của phân số $\dfrac{4}{9}$:
\choiceTF
{\True Tử số là $4$, mẫu số là $9$.}
{\True Phân số này đọc là Bốn phần chín.}
{Mẫu số chỉ số phần bằng nhau được lấy đi.}
{\True Mẫu số chỉ tổng số phần bằng nhau được chia ra.}
\loigiai{
\begin{itemize}
    \item Tử số là 4, mẫu số là 9: Đúng.
    \item Đọc là Bốn phần chín: Đúng.
    \item Tử số mới là số phần lấy đi (không phải mẫu số). Sai.
    \item Mẫu số chỉ tổng số phần chia đều: Đúng.
\end{itemize}
\begin{itemize}
    \item a) \textbf{Đúng}.
    \item b) \textbf{Đúng}.
    \item c) \textbf{Sai}.
    \item d) \textbf{Đúng}.
\end{itemize}
}
\end{ex}

\begin{ex} %[4D10V53]
Một trang trại có 50 con gà, 30 con vịt và 20 con ngan. Xét tính đúng sai:
\choiceTF
{\True Tổng số gia cầm của trang trại là 100 con.}
{\True Phân số chỉ số gà là $\dfrac{50}{100}$.}
{Phân số chỉ số vịt là $\dfrac{30}{50}$.}
{\True Phân số chỉ số ngan là $\dfrac{20}{100}$.}
\loigiai{
\begin{itemize}
    \item Tổng: $50 + 30 + 20 = 100$ con. Đúng.
    \item Gà: $\dfrac{50}{100}$. Đúng.
    \item Vịt so với tổng: $\dfrac{30}{100} \neq \dfrac{30}{50}$. Sai.
    \item Ngan: $\dfrac{20}{100}$. Đúng.
\end{itemize}
\begin{itemize}
    \item a) \textbf{Đúng}.
    \item b) \textbf{Đúng}.
    \item c) \textbf{Sai}.
    \item d) \textbf{Đúng}.
\end{itemize}
}
\end{ex}

\textbf{Phần 3. Trả lời ngắn}
\setcounterpageref{ex}{0}

\begin{ex} %[4D10V53]
Tìm tử số của phân số $\dfrac{8}{15}$.
\shortans{8}
\loigiai{
\textbf{Đáp số:} $8$\\
Tử số là 8.
}
\end{ex}

\begin{ex} %[4D10V53]
Tìm mẫu số của phân số $\dfrac{11}{24}$.
\shortans{24}
\loigiai{
\textbf{Đáp số:} $24$\\
Mẫu số là 24.
}
\end{ex}

\begin{ex} %[4D10V53]
Một tuần lễ có 7 ngày. Phân số chỉ số ngày đi học trong tuần (từ thứ Hai đến thứ Sáu - 5 ngày) là bao nhiêu (viết dạng a/b)?
\shortans{5/7}
\loigiai{
\textbf{Đáp số:} $5/7$\\
$5/7$.
}
\end{ex}

\begin{ex} %[4D10V53]
Một hộp có 20 viên bi, trong đó có 9 viên bi đỏ, còn lại là bi xanh. Viết phân số chỉ số bi xanh so với tổng số bi (dạng a/b).
\shortans{11/20}
\loigiai{
\textbf{Đáp số:} $11/20$\\
Bi xanh: $20 - 9 = 11 \Rightarrow 11/20$.
}
\end{ex}

\begin{ex} %[4D10V53]
Viết phân số có tử số là số lẻ lớn nhất có 1 chữ số và mẫu số là 10 (dạng a/b).
\shortans{9/10}
\loigiai{
\textbf{Đáp số:} $9/10$\\
$9/10$.
}
\end{ex}

\begin{ex} %[4D10V53]
Có bao nhiêu phân số có tổng của tử số và mẫu số bằng 5 (với tử số và mẫu số là các số tự nhiên khác 0)?
\shortans{4}
\loigiai{
\textbf{Đáp số:} $4$\\
Các cặp (tử, mẫu): (1,4), (2,3), (3,2), (4,1) $\Rightarrow 4$ phân số.
}
\end{ex}
